\noindent I set out to ask which mechanism produces homophily in the networks of newcomers. The short answer is that  the same-origin composition of these workers' networks is inherited rather than produced inside the Park. This is expected given the nature of my data, as we see the beginning of network formation (and nothing else). Within HIP, from what we can see, ties seem to depend on prolonged exposure from sharing a unit and whether the worker can be understood. \\


\noindent \textit{Importation.} This is the channel that carries the same-origin premium. Two workers from the same woreda are \pnChanImportOriginRel{} times as likely as any other pair to have arrived already acquainted. The premium is hyperlocal, \textit{i.e.} depends on woreda ties rather than zones. A shared woreda raises the chance that a worker produces a co-hire's name at all, but once the pairs she arrived knowing are set aside, that premium falls to the low end of the bracket reported in Section~\ref{sec:results}. \\

\noindent \textit{Opportunity.} Among the pairs who arrived as strangers, the firm's having placed the two in the same production unit multiplies the chance of a tie by \pnChanOppUnitRel{}. Two workers who cannot converse form a tie at \pnChanOppKappa{} of the rate of two who can. Both are estimated on pairs who met for the first time at HIP. \\

\noindent \textit{Taste.} At the woreda level I find nothing, and it is probably because I cannot detect it. The point estimate is \pnChanTasteTauPp{} of a percentage point, but the smallest effect this design would catch is \pnChanTasteMdePp{} percentage points (calculating the minimum detectable effect), which is the base tie rate itself. I can however rule out a big taste effect. On depth, shared origin does not make a tie deeper — it makes it slightly thinner, significantly so on both rosters once the worker is held fixed (\pnMuxOriginCohire{} on the co-hire side, \pnMuxOriginGn{} in the general network). Holding the layer profiles of kinship, vintage and co-residence, most of what origin predicts about tie content disappears. \\

\noindent \textit{Exclusion.} Among majority respondant the penalty insignificant. Minority workers name each other markedly more. This is also underpowered.

\noindent  If these results held with a better identification strategy and possibly a panel, this may imply that the institution has actually a big impact on how homophilous the newcomers' network can be. How it recruits and how it assigns production units can shape these people's networks. \\

\parhead{Limitations}

\noindent The first limitation is that this is a single cross-section, taken once, a median of \pnHireLagMedian{} days after the workers were hired. A panel would replace the whole construction, as I could observe the evolutions of ties, and assess if kinship becomes less prevalent over time. From my understanding, kin represents the key social and insurance institution in Ethiopia, so this would be a very important aspect to study. \\

\noindent The second limitation is power, and it falls on the two channels I do not detect. Taste rests on \pnChanTasteOriginPairs{} same-woreda pairs who arrived as strangers, of whom \pnChanTasteTies{} went on to be tied. Exclusion rests on the pairs in which exactly one of the two is a minority worker, \pnChanExclOutPairs{} of them. Nulls estimated on that much variation are not evidence of absence and are not reported as such anywhere in this thesis. \\
